\chapter{EVOLUÇÃO METODOLÓGICA DO PROJETO}
\label{ch:evolucao}

Este capítulo reconstrói o percurso metodológico que levou o projeto de uma avaliação binária de notícias reescritas à análise multidimensional de deriva em cadeias. O objetivo não é antecipar os resultados do estudo, apresentados no Capítulo \ref{ch:resultados}, mas explicitar os problemas encontrados, as decisões tomadas e as limitações que motivaram cada mudança de desenho. Parte das análises iniciais teve caráter exploratório e empregou configurações diferentes daquela prevista para a etapa consolidada; por isso, elas são apresentadas como evidências do processo de definição do método, e não como resultados finais.

A evolução ocorreu em quatro movimentos. Primeiro, foi construído um pipeline capaz de coletar notícias, reescrevê-las sob perspectivas predefinidas e registrar as saídas de diferentes provedores. Em seguida, sucessivas rodadas de classificação e auditoria mostraram que rótulos binários não descreviam satisfatoriamente as transformações produzidas. A terceira etapa introduziu o \textit{Structured Topic Drift Index} (STDI) e ampliou progressivamente sua representação de conteúdo e enquadramento. Por fim, o desenho foi organizado em cadeias controladas, nas quais as mesmas notícias percorrem sequências distintas de perspectivas. As especificações vigentes de dados, modelos, fórmulas e análises são apresentadas no Capítulo \ref{ch:metodos}; neste capítulo, versões anteriores são discutidas apenas quando ajudam a explicar uma decisão metodológica.

\section{PIPELINE INICIAL DE REESCRITA E AUDITORIA}
\label{sec:pipeline-inicial}

O desenvolvimento começou com um pipeline para obtenção de notícias por uma interface de programação, armazenamento tabular, reescrita por LLM e auditoria das versões geradas. A coleta inicial reuniu textos em inglês e preservou, quando disponíveis, campos como título, descrição, conteúdo, idioma, país, categoria, fonte, data de publicação e identificador do artigo. A finalidade dessa etapa era estabelecer uma infraestrutura rastreável: para cada execução, o sistema deveria registrar a entrada utilizada, o provedor, o modelo, a perspectiva, o \textit{prompt}, a saída, o estado da requisição e eventuais erros.

Antes das execuções em maior escala, foram realizados testes de funcionamento com amostras reduzidas. Esses testes verificaram a coleta, a seleção da coluna textual, a comunicação com diferentes provedores, a persistência das reescritas e a execução sequencial dos notebooks de auditoria. O histórico do projeto mostra que o pipeline foi ampliado de forma incremental: primeiro foram adicionados os geradores, depois a coleta de notícias, o detector pré-treinado, a organização das execuções em diretórios próprios e os relatórios automáticos. Essa infraestrutura permitiu preservar não apenas resultados agregados, mas também as saídas de cada modelo e os notebooks executados.

As primeiras condições sistemáticas de reescrita foram definidas por quatro perspectivas: direita conservadora, esquerda progressista, negacionista conspiratória e cética investigativa. Cada perspectiva era expressa por uma instrução textual que orientava enquadramento, ênfase, tom e vocabulário. O \textit{prompt} também solicitava a preservação do conteúdo factual e proibia a invenção de dados, números, citações, acontecimentos ou personagens. Essas instruções constituíam condições computacionais de geração; não se assumia que correspondessem a personalidades humanas estáveis nem que fossem seguidas com a mesma intensidade por todos os modelos.

A matriz principal combinou as quatro perspectivas com três conjuntos de notícias: um conjunto geral, um conjunto de política e um conjunto de ciência. Em cada combinação, foram solicitadas até 50 reescritas a quatro configurações: Gemini 3.1 Flash Lite Preview, NVIDIA Nemotron 3 Super, StepFun 3.5 Flash e Llama 3.1 8B. As taxas de sucesso variaram em razão de recusas, indisponibilidade e respostas inválidas, mas cada execução preservou o denominador e os erros. Em uma rodada adicional sobre notícias de política, a cobertura de modelos foi ampliada com GPT-4.1 Mini e Grok 4.1 Fast Non-Reasoning.

O pipeline executava duas formas de auditoria. A primeira aplicava o classificador pré-treinado \texttt{jy46604790/Fake-News-Bert-Detect} diretamente ao texto reescrito. A segunda comparava a notícia original e a reescrita com \texttt{MoritzLaurer/mDeBERTa-v3-base-mnli-xnli}. Nessa auditoria pareada, uma versão era sinalizada quando a probabilidade de contradição atingia o limiar de 0,60 e superava a probabilidade de implicação por pelo menos 0,20. Os procedimentos não tinham o mesmo alvo: o primeiro atribuía um rótulo ao documento reescrito, enquanto o segundo estimava a relação textual entre duas versões.

Essa primeira arquitetura foi útil para testar a geração em diferentes domínios, perspectivas e modelos, além de revelar uma limitação conceitual. Mesmo com os relatórios detalhados, uma saída binária não indicava qual parte da notícia havia mudado nem separava alterações de conteúdo, contradições entre versões e mudanças de enquadramento. O passo seguinte não abandonou os testes anteriores; procurou construir uma medida compatível com aquilo que eles não conseguiam representar.

\section{LIMITES DA CLASSIFICAÇÃO BINÁRIA}
\label{sec:limites-classificacao-binaria}

A decisão de desenvolver uma medida de deriva não resultou de um único teste. A matriz principal correspondeu ao cruzamento completo entre três conjuntos de notícias --- geral, política e ciência --- e quatro perspectivas --- direita conservadora, esquerda progressista, negacionista conspiratória e cética investigativa. Assim, foram consideradas as combinações geral--conservadora, geral--progressista, geral--conspiratória e geral--cética; política--conservadora, política--progressista, política--conspiratória e política--cética; e ciência--conservadora, ciência--progressista, ciência--conspiratória e ciência--cética. Como cada combinação previa 50 notícias reescritas por quatro modelos, o desenho totalizava 2.400 solicitações de reescrita. As execuções selecionadas para compor essa matriz produziram 2.327 saídas válidas, todas submetidas ao detector pré-treinado. Nenhuma recebeu a classe falsa. A auditoria pareada conseguiu processar 2.281 pares e sinalizou 51 como potencialmente falsos após a reescrita, segundo a regra de contradição descrita na Seção \ref{sec:pipeline-inicial}. Portanto, o padrão foi observado em diferentes conjuntos, perspectivas e modelos de geração, e não apenas em uma amostra isolada.

A ampliação posterior com GPT-4.1 Mini e Grok 4.1 Fast Non-Reasoning acrescentou 400 reescritas bem-sucedidas sobre notícias de política. O detector pré-treinado novamente não atribuiu a classe falsa a nenhuma saída. A auditoria pareada processou 392 pares e sinalizou cinco. Essa rodada reforçou o padrão anterior com dois geradores adicionais e mostrou que a diferença entre os dois procedimentos não dependia somente das quatro configurações usadas na matriz principal.

Essas rodadas cumpriram uma função diagnóstica bem definida. Elas fizeram variar o domínio das notícias, a perspectiva solicitada e o modelo gerador, mantiveram os textos produzidos e registraram falhas de execução. Não foram planejadas como um \textit{benchmark} de desempenho de detectores, pois as reescritas não receberam rótulos factuais humanos e o próprio \textit{prompt} solicitava a preservação dos fatos. O objetivo era verificar se as saídas dos auditores eram suficientes para estudar a transformação causada pela reescrita.

O resultado desse diagnóstico foi a identificação de duas limitações complementares. O classificador pré-treinado produziu uma resposta estável, porém não localizou mudanças graduais entre versões. A auditoria de inferência textual identificou um subconjunto de relações contraditórias, mas reduziu cada par a uma decisão baseada em limiares. Além disso, uma contradição entre textos não estabelece qual deles corresponde melhor a evidências externas. Como discutido na Seção \ref{sec:limites-deteccao}, classificadores também podem ser sensíveis à autoria e ao estilo de geração \cite{su2024adapting,wu2024sheepdog}.

Assim, a mudança de método foi motivada pelo acúmulo de resultados consistentes com a inadequação da representação binária para a pergunta do projeto. Em vez de perguntar apenas qual rótulo seria atribuído à reescrita, tornou-se necessário medir quanto e de que maneira ela se diferenciava da notícia de origem. A nova formulação exigia uma unidade de análise pareada e uma medida decomponível, sem converter mudança textual em julgamento automático de verdade.

\section{DESENVOLVIMENTO DO STDI}
\label{sec:desenvolvimento-stdi}

A primeira implementação do STDI introduziu uma representação estruturada para cada texto, composta por tema principal, subtópicos, entidades centrais, relações no formato sujeito--ação--objeto e enquadramento narrativo opcional. A comparação transformava essas estruturas em componentes interpretáveis de mudança. Desde essa etapa, a análise de cadeias distinguia a distância de cada versão ao texto original, a mudança em relação à versão imediatamente anterior e a soma das mudanças incrementais.

Para a formulação adotada no TCC, os quatro componentes de conteúdo --- tema, subtópicos, entidades e relações --- são tratados simetricamente, com pesos iguais. Versões intermediárias do código experimentaram ponderações diferentes, e alguns arquivos de comparação preservados em \texttt{output/topic\_drift/} foram produzidos antes da equalização. Como as simulações finais serão refeitas, esses escores intermediários não serão transportados diretamente para a análise final. Caso algum desses pares seja reutilizado, o índice deverá ser recalculado com a fórmula consolidada; as ponderações históricas não serão usadas para justificar resultados posteriores.

A representação estruturada resolveu parte do problema de interpretabilidade, pois permitia localizar a dimensão responsável por um escore. Entretanto, a comparação literal de rótulos e conjuntos extraídos permanecia sensível a paráfrases. Duas expressões semanticamente próximas podiam ser tratadas como diferentes, enquanto um mesmo rótulo poderia ocultar mudanças no modo de apresentação. A extração por LLM também introduzia uma segunda fonte de variação: antes de comparar os textos, o sistema dependia da estabilidade do modelo responsável por convertê-los em estruturas.

O projeto incorporou dois sinais complementares. O primeiro foi a contradição interna, inicialmente binária e depois representada por um valor contínuo de severidade e centralidade. O segundo foi a variação em valência, ativação e dominância (VAD), acrescentada para representar mudanças afetivas que não eram capturadas pelos componentes de conteúdo. Conforme discutido na Seção \ref{sec:representacao-afetiva}, VAD descreve propriedades estimadas do texto e não o estado emocional de autores ou leitores \cite{buechel2017emobank,mohammad2018vad}. Do mesmo modo, contradição interna não equivale a falsidade factual.

As análises exploratórias de VAD compararam distribuições de notícias rotuladas, pares aproximados por tema e contextos sintéticos de polaridade diferente. Esses ensaios mostraram que o sinal era sensível ao idioma, ao modelo e ao contexto da entrada. Por essa razão, VAD deixou de ser interpretado como marcador isolado de notícia falsa e passou a compor a descrição do enquadramento afetivo. A decisão metodológica foi preservar seus três componentes separadamente e limitar sua contribuição ao índice combinado.

Em uma etapa de calibração, foi preparado um conjunto de 50 pares controlados. Cada par associava uma notícia original a uma reescrita produzida para enfatizar uma dimensão específica, como tema, subtópico, entidade, relação, contradição ou enquadramento afetivo. A mesma extração estrutural foi reutilizada em dois caminhos de comparação. Um deles solicitava a um LLM uma avaliação semântica dos componentes; o outro utilizava embeddings do modelo \texttt{sentence-transformers/all-MiniLM-L6-v2}, agrupamento e regras fixas. Essa organização permitiu comparar métodos sem repetir a extração e confundir diferenças de representação com diferenças de comparação.

Para construir os grupos de referência usados nas regressões, o projeto partiu de um conjunto de notícias previamente rotuladas como falsas. Quinhentas dessas notícias foram reescritas pelo modelo \texttt{gpt-5.6-luna} mediante um processo de verdadeirização: produzir uma versão neutra sobre o mesmo tema, preservar aproximadamente o estilo, a estrutura e o comprimento do texto e corrigir ou remover alegações não sustentadas, sem inventar fontes, citações, números ou acontecimentos. Cada original rotulado como falso foi pareado com sua versão reescrita; os dois grupos foram então descritos por atributos do STDI, propriedades do texto e representações TF--IDF.

Também foram executadas auditorias de cobertura temática e regressões sobre esses grupos de referência. Nessa execução, não foi incluído um conjunto independente de notícias verdadeiras, e as versões verdadeirizadas não foram verificadas contra fontes externas. Portanto, os rótulos usados na regressão distinguem originais rotulados como falsos de reescritas solicitadas como verdadeiras, e não classes factuais confirmadas. A finalidade dessas análises era diagnóstica: identificar variáveis associadas aos grupos e procurar atalhos lexicais ou estilísticos. Um bom desempenho preditivo nesse contexto não transforma o STDI em verificador factual, porque o procedimento de reescrita e os sinais de autoria podem explicar parte da separação.

A fórmula foi então organizada hierarquicamente. Os componentes de conteúdo formam a base do índice; contradição e VAD entram como incrementos limitados pela distância ainda disponível até o valor máximo. Essa composição evita que a soma ultrapasse a escala do índice e mantém visível a contribuição de cada dimensão. A fórmula consolidada e as análises de sensibilidade são apresentadas no Capítulo \ref{ch:metodos}.

\section{CONSOLIDAÇÃO DAS CADEIAS DE REESCRITA}
\label{sec:consolidacao-cadeias}

Em paralelo ao desenvolvimento do índice, foi criada uma simulação em grafo para encadear reescritas. Cada nó recebe o texto produzido pelo nó anterior, aplica sua perspectiva e gera a entrada da etapa seguinte. Essa estrutura permite observar a trajetória da informação, e não somente pares independentes. As primeiras execuções do grafo verificaram a passagem de estado, a ordem dos nós, a persistência das saídas e a integração das métricas; elas não constituem a execução definitiva do experimento.

O desenho atual prevê 16 cenários, cada um com quatro posições, e um vocabulário de seis perspectivas: direita conservadora (C), esquerda progressista (P), negacionista conspiratória (D), cética investigativa (S), amplificadora emocional (E) e comunicadora conciliatória (M). As duas últimas foram acrescentadas para testar, respectivamente, intensificação afetiva e mediação entre enquadramentos. As letras identificam instruções de geração, não categorias psicológicas observadas em pessoas.

Os cenários foram organizados para permitir contrastes interpretáveis. A sequência SSSS funciona como referência de reescrita cética repetida. CCCC, PPPP e DDDD formam cadeias homogêneas. CCPP e PPCC invertem a ordem de dois blocos ideológicos, enquanto CPCP e PCPC alternam essas perspectivas. DDSS e SSDD invertem blocos conspiratórios e céticos; DSDS e SDSD alternam as mesmas duas condições. DDES e DDSE deslocam a amplificação emocional para antes ou depois da perspectiva cética. Por fim, CMPC e CPMC posicionam a comunicação conciliatória entre perspectivas ou depois do encontro entre elas. Esses pares permitem examinar associação com a ordem mantendo constante, sempre que possível, a composição da cadeia.

A execução preliminar aplicou os cenários às mesmas 50 notícias com uma configuração comum de modelo. Esse arranjo permite que cada notícia funcione como seu próprio controle nas comparações entre sequências e reduz a variação decorrente da seleção de itens e do uso de modelos diferentes, presente nas etapas iniciais. Ele não elimina, porém, a estocasticidade da geração, a possível aderência desigual às perspectivas, a dependência do \textit{prompt} e a influência do comprimento ou do domínio da notícia.

As simulações serão executadas novamente após o congelamento da configuração final de corpus, modelo, \textit{prompt}, parâmetros e tratamento de falhas. Por isso, os totais e escores atualmente armazenados são tratados apenas como diagnóstico do pipeline e não serão apresentados como resultados definitivos. Após a nova execução, esta seção deverá ser conferida para refletir somente o desenho efetivamente utilizado; os valores finais pertencerão ao Capítulo \ref{ch:resultados}.

Para cada etapa válida, o sistema preserva a notícia original, a versão recebida, a reescrita, a perspectiva, o modelo, os estados de processamento, a estrutura extraída, os sinais VAD e os componentes do índice. São calculadas três visões da trajetória: STDI em relação ao original, STDI incremental em relação ao passo anterior e STDI cumulativo como soma dos incrementos. O primeiro resume o afastamento da origem; o segundo localiza mudanças entre posições adjacentes; o terceiro registra a quantidade total de transformação percorrida, ainda que uma etapa posterior se reaproxime do texto original.

Esse desenho encerra a passagem de uma pergunta binária para uma análise pareada e longitudinal. O foco empírico deixa de ser o rótulo atribuído a uma saída isolada e passa a ser a trajetória de componentes informacionais e afetivos sob sequências controladas. As comparações entre cenários são formuladas como associações observadas nas saídas do modelo. Atribuições causais a uma perspectiva ou generalizações para comportamento humano exigiriam controles e evidências que não fazem parte desta simulação.

Em síntese, a evolução metodológica produziu três separações que estruturam o restante do trabalho: entre deriva e veracidade, entre mudança incremental e distância ao original, e entre a geração textual e sua mensuração posterior. O Capítulo \ref{ch:metodos} detalha a configuração consolidada que implementa essas decisões; o Capítulo \ref{ch:resultados} avaliará seu comportamento após a nova execução das cadeias.
